%! AP Statistics — Unit 2, Lesson 4
\documentclass[10pt]{article}
\usepackage[margin=0.6in]{geometry}
\usepackage{xcolor}
\usepackage[most]{tcolorbox}
\usepackage{enumitem}
\usepackage{multicol}
\usepackage{amsmath,amssymb}
\usepackage{array}
\usepackage{tabularx}
\tcbuselibrary{breakable}

\definecolor{sky}{HTML}{EAF3FB}
\definecolor{navy}{HTML}{1F3A5F}
\definecolor{redbox}{HTML}{FCECEC}
\definecolor{greenbox}{HTML}{EDF8F0}
\definecolor{goldbox}{HTML}{FFF7E6}
\newtcolorbox{skillbox}[2][]{
    colback=#2, colframe=navy, boxrule=0.8pt, arc=2mm,
    left=2mm,right=2mm,top=1.5mm,bottom=1.5mm,
    fonttitle=\bfseries, title={#1}, halign title=center, breakable
}
\setlist[itemize]{leftmargin=1.2em,itemsep=0.35em,topsep=0.35em}
\setlist[enumerate]{leftmargin=1.4em,itemsep=0.35em,topsep=0.35em}
\newcolumntype{C}[1]{>{\centering\arraybackslash}p{#1}}
\newcolumntype{Y}{>{\centering\arraybackslash}X}

\newcommand{\CourseName}{AP Statistics}
\newcommand{\SchoolYear}{2026--2027}
\newcommand{\MeetingLength}{55 minutes}
\newcommand{\UnitNumberName}{Unit 2: Exploring Two-Variable Data \quad}
\newcommand{\LessonNumberName}{Lesson 2.4: Representing the Relationship Between Two Quantitative Variables}

\begin{document}
\begin{center}
    {\Large\bfseries \CourseName: \SchoolYear\ (\MeetingLength) \\
    {\small\bfseries \UnitNumberName \LessonNumberName}}
\end{center}

\begin{tcolorbox}[colback=sky,colframe=navy,boxrule=0.9pt,arc=2mm,
    left=2mm,right=2mm,top=1.5mm,bottom=1.5mm]
    \textbf{Primary Objective:} Students will construct and interpret scatterplots for two
    quantitative variables. Students will describe the relationship using the DOFS framework:
    \textbf{D}irection, \textbf{F}orm, \textbf{S}trength, and \textbf{O}utliers, in context
    (Big Idea: UNC, DAT; AP Skills 2 \& 4).
\end{tcolorbox}

\begin{skillbox}[Priority Ideas \& Skills]{goldbox}
\begin{minipage}[t]{0.48\linewidth}
    \textbf{AP Skills 2 \& 4 --- Data Analysis \& Statistical Argumentation}
    \begin{itemize}
        \item Construct a correctly labeled scatterplot with the explanatory variable on the $x$-axis.
        \item Describe the relationship using direction (positive/negative/none), form (linear/curved), strength (strong/moderate/weak), and outliers.
        \item Write a complete description in context.
    \end{itemize}
\end{minipage}
\hfill
\begin{minipage}[t]{0.48\linewidth}
    \textbf{Key Understandings}
    \begin{itemize}
        \item A scatterplot shows the relationship between two quantitative variables for the same set of individuals.
        \item ``Positive association'' means as $x$ increases, $y$ tends to increase (not that the relationship is good).
        \item The form of the relationship (linear vs.\ curved) determines which model is appropriate.
        \item An outlier in a scatterplot is a point that does not follow the overall pattern.
    \end{itemize}
\end{minipage}
\end{skillbox}

\begin{skillbox}[{Vocabulary, Concepts, \& Theorems}]{greenbox}
\begin{minipage}[t]{0.48\linewidth}
    \begin{itemize}
        \item \textbf{Scatterplot} --- a graph of ordered pairs $(x, y)$ for two quantitative variables measured on the same individual
        \item \textbf{Positive association} --- as $x$ increases, $y$ tends to increase
        \item \textbf{Negative association} --- as $x$ increases, $y$ tends to decrease
        \item \textbf{No association} --- knowing $x$ gives no information about $y$
    \end{itemize}
\end{minipage}
\hfill
\begin{minipage}[t]{0.48\linewidth}
    \begin{itemize}
        \item \textbf{Form} --- the overall shape of the pattern: linear, curved, or none
        \item \textbf{Strength} --- how closely the points follow the form; strong, moderate, or weak
        \item \textbf{Outlier (bivariate)} --- a point that deviates from the overall pattern; may have extreme $x$, extreme $y$, or both
        \item \textbf{Influential point} --- an outlier whose removal would substantially change the line of best fit (preview for Lesson 2.8)
    \end{itemize}
\end{minipage}
\end{skillbox}

\begin{skillbox}[Activate Prior Knowledge \& Spiral Review]{sky}
\begin{minipage}[t]{0.48\linewidth}
    \textbf{Quick Review (3 min)}\\
    Recall: in a scatterplot from prior math courses, what goes on the $x$-axis vs.\ the $y$-axis?
    Connect to today: the \textit{explanatory} variable always goes on $x$, the \textit{response} on $y$.
\end{minipage}
\hfill
\begin{minipage}[t]{0.48\linewidth}
    \textbf{Spiral Connections}
    \begin{itemize}
        \item Lesson 2.1: explanatory vs.\ response variable determines axis placement.
        \item Lesson 2.5: correlation measures the strength and direction of a \textit{linear} association numerically.
        \item Lessons 2.6--2.9: linear regression fits a line to the scatterplot.
    \end{itemize}
\end{minipage}
\end{skillbox}

\begin{skillbox}[Hook \& Lesson]{sky}
\begin{minipage}[t]{0.48\linewidth}
    \textbf{Hook (5 min)}\\
    Display four unlabeled scatterplots: strong positive linear, weak negative linear, curved, and
    no association. Ask: ``Rank these from most useful for prediction to least useful.'' Debrief:
    the clearer the pattern, the better our predictions --- and today we learn the language to
    describe these patterns precisely.
\end{minipage}
\hfill
\begin{minipage}[t]{0.48\linewidth}
    \textbf{Lesson Flow (15 min)}
    \begin{enumerate}
        \item Construct a scatterplot from a small dataset (10 points) together on the board.
        \item Describe the plot using direction, form, strength, and outliers in complete sentences with context.
        \item Examine 3--4 additional scatterplots; students practice describing each.
        \item Discuss: what does it mean for a point to be an ``outlier'' in a bivariate context?
    \end{enumerate}
\end{minipage}
\end{skillbox}

\begin{skillbox}[Explicit Instruction \& Active Monitoring]{sky}
\begin{minipage}[t]{0.48\linewidth}
    \textbf{Direct Instruction (10 min)}
    \begin{itemize}
        \item Model response template: ``There is a [strong/moderate/weak] [positive/negative] [linear/curved] association between [explanatory variable] and [response variable] for these [individuals]. [Outlier sentence if applicable.]''
        \item Discuss that ``linear'' is an important qualifier --- correlation and regression (Lessons 2.5--2.8) only apply to linear relationships.
        \item Emphasize: always include context (name the variables and population).
    \end{itemize}
\end{minipage}
\hfill
\begin{minipage}[t]{0.48\linewidth}
    \textbf{Active Monitoring}
    \begin{itemize}
        \item Common errors: swapping $x$ and $y$ axes; saying ``the data is positive'' instead of naming the association.
        \item Prompt: ``Is the form \textit{linear}? Is the strength \textit{strong}? Both matter for what we do next.''
        \item Ask students to estimate: ``If $x = 50$, what would you predict for $y$?'' --- motivates the need for a model.
    \end{itemize}
\end{minipage}
\end{skillbox}

\begin{skillbox}[Group Work \& Differentiation]{redbox}
\begin{minipage}[t]{0.48\linewidth}
    \textbf{Group Activity (8--10 min)}\\
    Groups each receive a dataset (e.g., car age vs.\ price; temperature vs.\ ice cream sales;
    study time vs.\ test score). Groups:
    \begin{enumerate}
        \item Construct the scatterplot on graph paper.
        \item Write a full DOFS description.
        \item Identify any outliers and speculate on explanations.
        \item Share with the class and compare relationships across datasets.
    \end{enumerate}
\end{minipage}
\hfill
\begin{minipage}[t]{0.48\linewidth}
    \textbf{Differentiation}
    \begin{itemize}
        \item \textit{Support}: Provide pre-drawn axes; students only plot points and write the description.
        \item \textit{Extension}: Give a dataset where the relationship is curved (e.g., age vs.\ reaction time). Students explain why a straight line would be a poor model.
        \item \textit{Technology}: Plot using Desmos or a graphing calculator and compare to the hand-drawn version.
    \end{itemize}
\end{minipage}
\end{skillbox}

\begin{skillbox}[Individual Work \& Assessment]{redbox}
\begin{minipage}[t]{0.48\linewidth}
    \textbf{Exit Ticket (5 min)}\\
    Display a scatterplot of hours of TV per day vs.\ GPA for 25 students. Students write a
    complete DOFS description in 3--4 sentences, including context and any notable outliers.
\end{minipage}
\hfill
\begin{minipage}[t]{0.48\linewidth}
    \textbf{AP-Style Check}\\
    \textit{``Describe the relationship between [x-variable] and [y-variable] shown in the scatterplot.''}\\[4pt]
    Model: ``There is a moderate negative linear association between daily TV hours and GPA for
    these 25 students. As TV hours increase, GPA tends to decrease. One student who watches
    7 hours daily but has a 3.8 GPA appears to be an outlier.''
\end{minipage}
\end{skillbox}

\begin{skillbox}[Reinforcement \& Extension]{goldbox}
\begin{minipage}[t]{0.48\linewidth}
    \textbf{Homework / Practice}
    \begin{itemize}
        \item Construct and describe scatterplots for 2 provided datasets.
        \item Match 4 DOFS descriptions to 4 scatterplot images.
    \end{itemize}
\end{minipage}
\hfill
\begin{minipage}[t]{0.48\linewidth}
    \textbf{Extension / Preview}
    \begin{itemize}
        \item \textit{Extension}: Find a published scatterplot in a scientific article or news report.
              Write a DOFS description and assess whether the authors' conclusions are supported by the graph.
        \item \textit{Preview}: Lesson 2.5 quantifies the strength and direction of a \textit{linear}
              association with a single number --- the correlation coefficient $r$. Predict: what value
              would $r$ have for each of today's scatterplots?
    \end{itemize}
\end{minipage}
\end{skillbox}

\end{document}
